% !TeX root = elegantnote-en.tex
\section{Amines}

\subsection{Introduction 20260421}

\begin{enumerate}
    \item Amines are organic derivatives of ammonia with one or more alkyl or aryl groups bonded to the nitrogen atom.
    \item dopamine, %多巴胺
    epinephrine, %肾上腺素
    L-tryptophan, %色氨酸
    piperazine, %哌嗪
    \item \textit{Alkaloids}
\end{enumerate}

\subsection{Nomenclature of Amines 20260421 20260423}

\begin{enumerate}
    \item Amines are classified as \textbf{primary} (1°), \textbf{secondary} (2°), or \textbf{tertiary} (3°), quaternary (4°) corresponding to one, two, three or four alkyl or aryl groups bonded to nitrogen.
    \item \textbf{Quaternary ammonium salts} have four alkyl or aryl bonds to a nitrogen atom.
\end{enumerate}

\subsubsection{Common Names}

\begin{enumerate}
    \item Common names of amines are formed from the names of the alkyl groups bonded to nitrogen, followed by the suffix \textit{-amine}.
    \item In naming amines with more complicated structures, the \ce{-NH2} group is called the amino group. Secondary and tertiary amines are named by classifying the nitrogen atom (together with its alkyl groups) as an alkylamino group.
    \item Phenylamine is called \textit{aniline}, for example, and its derivatives are named as derivatives of aniline.
    \item Nitrogen heterocycles. The heteroatom is usually assigned position number 1. aziridine, pyrrole, pyrrolidine, imidazole, indole, pyridine, piperidine, pyrimidine, purine%pic
\end{enumerate}

\subsubsection{IUPAC Names}

\begin{enumerate}
    \item The IUPAC nomenclature for amines is similar to that for alcohols.
    \item The longest continuous chain of carbon atoms determines the root name. The \textit{-e} ending in the alkane name is changed to \textit{-amine}, and a number shows the position of the amino group along the chain.
    \item butan-2-amine, \textit{N}-methylbutan-2-amine, 2,4,\textit{N},\textit{N}-tetramethylhexan-3-amine
\end{enumerate}

\subsection{Structure of Amines 20260423}

\begin{enumerate}
    \item Ammonia has a slightly distorted tetrahedral shape.
    \item A tetrahedral amine with three different substituents (and a lone pair) is nonsuperimposable on its mirror image, and appears to be a chirality center.
    \item Cannot be resolved two enantiomers because the enantiomers interconvert rapidly, \textbf{nitrogen inversion}. %pic
    \begin{figure}[h]
        \centering
        \includegraphics[width=0.8\textwidth]{19.3_Nitrogen_Inversion.png}
        \caption{Nitrogen Inversion}
        \label{19.3_Nitrogen_Inversion}
    \end{figure}
    \item Several types of chiral amines can be resolved:
    \begin{enumerate}
        \item Amines whose chirality stems from the presence of asymmetric carbon atoms. %pic
        \item Quaternary ammonium salts with asymmetric nitrogen atoms(larger than hydrogen). %pic
        \item Amines that cannot attain the \textit{sp}$^2$ transition state for nitrogen inversion.
    \end{enumerate}
\end{enumerate}

\begin{example} % 19-4
    Which of the amines listed next can be resolved into enantiomers? In each case, explain why interconversion of the enantiomers would or would not take place. 
    \begin{enumerate}[label=(\alph*)]
        \item \textit{cis}-2-methylcyclohexanamine
        \item \textit{N}-ethyl-\textit{N}-methylcyclohexanamine
        \item \textit{N}-methylaziridine
        \item ethylmethylanilinium iodide
        \item methylethylpropylisopropylammonium iodide
    \end{enumerate}
\end{example}

\begin{example} % 19-34
    Which of the following compounds are capable of being resolved into enantiomers?
\end{example}

\subsection{Physical Properties of Amines 20260423}

\begin{enumerate}
    \item Nitrogen is less electronegative than oxygen. Amines form weaker hydrogen bonds than do alcohols of similar molecular weights.
    \item Amines tend to be soluble in alcohols, and the lowermolecular-weight amines (up to about four carbon atoms) are relatively soluble in water.
    \item Perhaps the most obvious property of amines is their characteristic odor of rotting fish.
    putrescine， %腐胺 
    cadaverine %尸胺
\end{enumerate}

\subsection{Basicity of Amines 20260423}

\begin{enumerate}
    \item An amine is a nucleophile (a Lewis base) because its lone pair of nonbonding electrons can form a bond with an electrophile. An amine can also act as a Brønsted–Lowry base by accepting a proton from a proton acid.
    \item \textbf{Base-dissociation constant} for the amine, symbolized by $\textbf{K}_{\mathrm{b}}$.
\end{enumerate}

\subsection{Effects on Amine Basicity 20260423 20260428}

\begin{enumerate}
    \item Any structural feature that stabilizes the ammonium ion (relative to the free amine) shifts the reaction toward the right, making the amine a stronger base.
    \item Substitution by Alkyl Groups
    \begin{enumerate}
        \item Alkyl groups are electron-donating toward cations.
        \item The simple alkylamines are generally stronger bases than ammonia.
    \end{enumerate}
    \item Resonance Effects on Basicity
    \begin{enumerate}
        \item Arylamines (anilines ($\mathrm{p}K_{\mathrm{b}} = 9.40$) and their derivatives) are much weaker bases than simple aliphatic amines. This reduced basicity is due to resonance delocalization of the nonbonding electrons in the free amine.
        \item Resonance effects also influence the basicity of pyrrole. When the pyrrole nitrogen is protonated, pyrrole loses its aromatic stabilization.
    \end{enumerate}
    \item Hybridization Effects
    \begin{enumerate}
        \item Electrons are held more tightly by orbitals with more s character.
        \item Pyridine’s nonbonding electrons are less available for bonding to a proton. ($\mathrm{p}K_{\mathrm{b}} = 8.75$)
    \end{enumerate}
\end{enumerate}

\begin{example} % 19-6
    Rank each set of compounds in order of increasing basicity.
    \begin{enumerate}[label=(\alph*)]
        \item \ce{NaOH}, \ce{NH3}, \ce{CH3NH2}, \ce{Ph-NH2}
    \end{enumerate}
\end{example}

\begin{example} % 19-33
    Rank the amines in each set in order of increasing basicity.
\end{example}

\begin{example} %19-35
    Complete the following proposed acid–base reactions, and predict whether the reactants or products are favored.
\end{example}

\subsection{Salts of Amines 2060428}

\begin{enumerate}
    \item Protonation of an amine gives an amine salt.
    \item Simple amine salts are named as the substituted \textbf{ammonium salts}. Salts of complex amines use the names of the amine and the acid that make up the salt.
    \item Amine salts are ionic, high-melting, nonvolatile solids, much more soluble in water than the parent amines, and only slightly soluble in nonpolar organic solvents.
    \item Formation of amine salts can be used to isolate and characterize amines.
    \item We can use the formation of amine salts to separate amines from less basic compounds. %pic
    \item Many drugs and other biologically important amines are commonly stored and used as their salts.
\end{enumerate}

\subsection{Phase Transfer Catalysts 20260428}

\subsection{Spectroscopy of Amines}

\subsubsection{Infrared Spectroscopy}

\begin{enumerate}
    \item The most reliable IR absorption of primary and secondary amines is the \chemfig{N-H} stretch whose frequency appears between \qty{3200}{} and \qty{3500}{\per\centi\meter}. 
\end{enumerate}

\subsubsection{Proton NMR Spectroscopy}

\subsubsection{Carbon NMR Spectroscopy}

\subsubsection{Mass Spectrometry}

\begin{enumerate}
    \item Nitrogen has an odd valence and an even mass number.
    \item Whenever an odd number of nitrogen atoms are present in a molecule, the molecular ion has an odd mass number. Most of the fragments have even mass numbers.
\end{enumerate}

\subsection{Reactions of Amines with Ketones and Aldehydes (Review) 20260428}

\begin{enumerate}
    \item 
\end{enumerate}

\subsection{Aromatic Substitution of Arylamines and Pyridine 20260428}

\subsubsection{Electrophilic Aromatic Substitution of Arylamines}

\begin{enumerate}
    \item The \ce{-NH3+} group is strongly deactivating (and meta-allowing). Therefore, strongly acidic reagents are unsuitable for substitution of anilines.
\end{enumerate}

\subsubsection{Electrophilic Aromatic Substitution of Pyridine}

\begin{enumerate}
    \item Pyridine resembles nitro benzene in its reactivity toward electrophilic aromatic substitution.
    \item Friedel–Crafts reactions fail completely, and other substitutions require unusually strong conditions.
    \item Deactivation results from the electron-withdrawing effect of the electronegative nitrogen atom.
\end{enumerate}

\subsubsection{Nucleophilic Aromatic Substitution of Pyridine}

\subsection{Alkylation of Amines by Alkyl Halides 20260428}

\begin{enumerate}
    \item Amines react with primary alkyl halides to give alkylated ammonium halides.
    \item Alkylation proceeds by the $\mathrm{S}_{\mathrm{N}}2$ mechanism, so it is not feasible with tertiary halides because they are too hindered.
    \item Even if just one equivalent of the halide is added, some amine molecules will react once, some will react twice, and some will react three times.
\end{enumerate}

\subsection{Acylation of Amines by Acid Chlorides 20260428}

\begin{enumerate}
    \item Primary and secondary amines react with acid halides to form amides.
    \item This reaction is a \textit{nucleophilic acyl substitution}: the replacement of a leaving group on a carbonyl carbon by a nucleophile.
\end{enumerate}

\begin{mechanism}
    Acylation of an Amine by an Acid Chloride
    \begin{itemize}
        \item Step 1: A nucleophile attacks the strongly electrophilic carbonyl group of the acid chloride to form a tetrahedral intermediate.
        \item Step 2: The tetrahedral intermediate expels chloride ion.
        \item Step 3: Loss of a proton gives the amide.
    \end{itemize}
\end{mechanism}

\begin{enumerate}[resume]
    \item The diminished basicity of amides can be used to advantage in \textbf{electrophilic aromatic substitutions}.
    \item If the amino group of aniline is acetylated to give acetanilide, the resulting amide is still activating and ortho, para-directing. %pic
    \item Aryl amino groups are frequently acylated before further substitutions are attempted on the ring, and the acyl group is removed later by acidic or basic hydrolysis
\end{enumerate}

\begin{example} % solved 19-1
    Show how you would accomplish the following synthetic conversion in good yield.

    \textbf{SOLUTION} \quad An attempted Friedel–Crafts acylation on aniline would likely meet with disaster. The free amino group would attack both the acid chloride and the Lewis acid catalyst.
\end{example}

\subsection{Formation of Sulfonamides 20264028}

\begin{enumerate}
    \item 
\end{enumerate}

\subsection{Amines as Leaving Groups: The Hofmann Elimination 20260428}

\begin{enumerate}
    \item An amine cannot undergo elimination directly because the leaving group would be an amide ion (\ce{^-NH2} or \ce{^-NHR}), which is a very strong base and a poor leaving group.
    \item An amino group can be converted to a good leaving group by exhaustive methylation, which converts it to a quaternary ammonium salt that can leave as a neutral amine.
    \item Exhaustive methylation is usually accomplished using methyl iodide.
    \item Elimination of the quaternary ammonium salt generally takes place by the E2 mechanism, which requires a strong base. To provide the base, the quaternary ammonium iodide is converted to the hydroxide salt by treatment with silver oxide.
    \item In the Hofmann elimination the product is commonly the \textit{least-substituted} alkene (\textit{Hofmann product}).
    \item Remember that the E2 mechanism requires an \textbf{anti-coplanar arrangement} of the proton and the leaving group.
    \item The Hofmann product predominates because elimination of one of the \chemfig{C1} protons involves a lower-energy, more probable transition state than the crowded transition state required for Zaitsev (\chemfig{C2-C3}) elimination.
    \item The least-substituted alkene is usually the favored product.
\end{enumerate}

\begin{example} % s 19-2
    Predict the major product(s) formed when the following amine is treated with excess iodomethane, followed by heating with silver oxide.
\end{example}

\begin{example} % 19-20
    Predict the major products formed when the following amines undergo exhaustive methylation, treatment with and heating.
    \begin{enumerate}[label=(\alph*)]
        \item hexan-2-amine
        \item 2-methylpiperidine
        \item \textit{N}-ethylpiperidine
    \end{enumerate}
\end{example}

\subsection{Oxidation of Amines; The Cope Elimination 20260430}

\begin{enumerate}
    \item sdf
\end{enumerate}

\begin{example} % 19-22
    MCPBA
\end{example}

\subsection{Reactions of Arenediazonium Salts 20260430 20260507}

\begin{enumerate}
    \item d
    \item These reactions of diazonium salts are \textbf{extremely useful} for solving aromatic synthesis problems. %pic
    \begin{figure}[h]
        \centering
        \includegraphics[width=0.8\textwidth]{19.17_Diazonium_Salts_Reactions.png}
        \caption{Reactions of Diazonium Salts}
        \label{19.17_Diazonium_Salts_Reactions}
    \end{figure}
    \item Replacement of the Diazonium Group by Hydroxide: Hydrolysis
    \begin{enumerate}
        \item Hydrolysis takes place when the acidic solution of an arenediazonium salt is warmed.
        \item This is a useful laboratory synthesis of phenols.
    \end{enumerate}
    \item Replacement of the Diazonium Group by Chloride, Bromide, and Cyanide: The Sandmeyer Reaction
    \begin{enumerate}
       \item The use of cuprous salts to replace arenediazonium groups is called the Sandmeyer reaction.
       \item The Sandmeyer reaction (using cuprous cyanide) is also an excellent method for attaching another carbon substituent to an aromatic ring.
    \end{enumerate}
    \item Replacement of the Diazonium Group by Fluoride and Iodide
    \begin{enumerate}
        \item When an arenediazonium salt is treated with fluoroboric acid \ce{HBF4}, the diazonium fluoroborate precipitates out of solution.
        \item Aryl iodides are formed by treating arenediazonium salts with potassium iodide.
    \end{enumerate}
    \item Reduction of the Diazonium Group to Hydrogen: Deamination of Anilines
    \item Diazonium Salts as Electrophiles: Diazo Coupling
    \begin{enumerate}
        \item Arenediazonium ions act as weak electrophiles in electrophilic aromatic substitutions.
        \item The product have the structure \chemfig{Ar-N=N-Ar'}.
    \end{enumerate}
    \item Summary: Reactions of Amines
\end{enumerate}

\subsection{Synthesis of Amines by Reductive Amination 20260507}

\begin{enumerate}
    \item \textbf{Reductive amination} is the most general amine synthesis, capable of adding a primary or secondary alkyl group to an amine.
    \item First we form an imine or oxime derivative of a ketone or aldehyde, and then reduce it to the amine.
    \item Primary Amines
    \begin{enumerate}
        \item Primary amines result from condensation of hydroxylamine (zero alkyl groups) with a ketone or an aldehyde, followed by reduction of the oxime.
        \item The oxime is reduced using catalytic reduction, lithium aluminum hydride, or zinc and HCl.
    \end{enumerate}
    \item Secondary Amines
    \begin{enumerate}
        \item Condensation of a primary amine with a ketone or aldehyde forms an N-substituted imine (a Schiff base).
        \item Reduction of the imine, using either \ce{LiAlH4} or \ce{NaBH4}, gives a secondary amine.
    \end{enumerate}
    \item Tertiary Amines
    \begin{enumerate}
        \item Condensation of a secondary amine with a ketone or aldehyde gives an iminium salt.
        \item Iminium salts are frequently unstable, so they are rarely isolated.
        \item Sodium triacetoxyborohydride [\ce{NaBH(OCOCH3)3} or \ce{NaBH(OAc)3}] is less reactive than sodium borohydride, and it does not reduce the carbonyl group.
    \end{enumerate}
\end{enumerate}

\subsection{Synthesis of  Amines by Acylation–Reduction 20260507}

\begin{enumerate}
    \item The second general synthesis of amines is \textbf{acylation–reduction}.
    \item Reduction of the amide by lithium aluminum hydride (\ce{LiAlH4}) gives the corresponding amine (no hydroxyl generated).
    \item Primary amines
    \item Secondary Amines
    \item Tertiary Amines
\end{enumerate}

\subsection{Syntheses Limited to Primary Amines 20260507}

\subsubsection{Direct Alkylation and Gabriel Synthesis}

\begin{enumerate}
    \item Phthalimide
    \item The phthalimide anion is a strong nucleophile, displacing a halide or tosylate ion from a good $\mathrm{S}_{\mathrm{N}}2$ substrate.
    \item Heating the \textit{N}-alkyl phthalimide with hydrazine displaces the primary amine, giving the very stable hydrazide of phthalimide.
\end{enumerate}

\subsubsection{Reduction of Azides and Nitriles}

\begin{enumerate}
    \item Formation and Reduction of Azides
    \item Formation and Reduction of Nitriles
    \begin{enumerate}
        \item cyanide ion is a good nucleophile
        \item Nitriles are reduced to primary amines by lithium aluminum hydride or by catalytic hydrogenation.
    \end{enumerate}
\end{enumerate}

\subsubsection{The Hofmann Rearrangement of Amides}

\begin{enumerate}
    \item In the presence of a strong base, primary amides react with chlorine or bromine to form shortened amines, with the loss of the carbonyl carbon atom.
    \item This reaction, called the Hofmann rearrangement, is used to synthesize primary and aryl amines.
    \item isocyanate, carbamic acid
    \item The Hofmann rearrangement, can produce primary amines with primary, secondary, tertiary alkyl group or aryl amines.
    \item The conformation will keep the same.
\end{enumerate}

\begin{example} % 7th ed. 19-40
    
\end{example}

\begin{example} % 7th ed. 19-34
    \textit{Curtius rearrangement}
\end{example}

Assigment 19-36(no classification), 41, 42

\begin{example} % 7th ed, 19-62
    Ketones and aldehydes react with primary amines to give imines.
\end{example}